
% try to also keep this to 1-1.5 pages 

%\section{Related Works} 
%\label{sec:RelatedWorks}
%
%This work falls into the realm of GPU kernel performance metric prediction using LLMs. 
%The closest neighboring realms are those of (a) using LLMs for parallel code generation/optimization, and (b) using LLMs for program analysis and execution simulation.
%We discuss existing works from our realm and its neighbors below.
%
%
%\textbf{LLMs for Parallel Code Generation.}
%
%Over the past few years, the utilization of LLMs for parallel code generation has exploded.
%The code optimization works can be largely broken up into two camps: (1) fine-tuning/training-based, and (2) out-of-the-box LLM-based.
%
%The fine-tuning works usually focus on creating a dataset of parallel codes 
%The dataset is often synthetic in nature, so extra efforts are made to 
%Given the natural language understanding of LLMs, and the complexity of \cite{talkadoshScopeAllYou2023}
%
%Talk about fine-tuned code LLMs (CodeBERT, Codex, Meta LLM Compiler, CodeGen, etc.
%
%\cite{ouyangKernelBenchCanLLMs2025}
%%In the past two years there have been multiple efforts to use LLMs for CUDA, OpenMP, OpenCL, and Triton code generation for GPUs.
%
%%\citet{ouyangKernelBenchCanLLMs2025} were one of the first that tasked LLMs to generate optimized CUDA kernels from PyTorch model architectures.
%%They demonstrated that profiler feedback significantly improves the generated codes' performance during refinement steps.
%%
%%\citet{nicholsCanLargeLanguage2024} is another work where
%%
%%\citet{wenMultiKernelBenchMultiPlatformBenchmark2025} 
%%\citet{liTritonBenchBenchmarkingLarge2025}
%%\citet{kaplanParaCodexProfilingGuidedAutonomous2026} is a more recent work, where 
%
%
%
%\textbf{RL for LLM-based GPU Code Generation.}
%
%Since the seminal DeepSeek-R1 paper \cite{} demonstrated the benefits of using reinforcement learning (RL) to fine-tune and distill LLM knowledge, a series of works sprung up using similar RL techniques to fine-tune LLMs for code-generation with an focus on execution speedup.
%
%\cite{dongKernelBlasterContinualCrossTask2026}
%\cite{nicholsIntegratingPerformanceTools2025}
%\cite{baronioKevinMultiTurnRL2025}
%%Another vein of work has been fine-tuning and creating custom LLM-based models for GPU code generation and optimization.
%
%%CUDA-L1
%%Kevin
%%Harshitha + Nichols work
%
%
%\textbf{Agentic Iterative/Feedback-Driven Code Optimization.}
%
%\cite{chenCUDALLMLLMsCan2025}
%\cite{}
%
%% High inference costs with agents
%% HPCToolKit + Hatchet paper (not sure what it was called)
%
%\textbf{LLMs for GPU Performance Prediction.}
%% Bolet
%% Omniwise
%
%
%\textbf{LLMs for Execution Simulation / Program Analysis / Trace Prediction.}
%
%\cite{lyuLargeLanguageModels2024} simulated code execution of Python, Java, C, C++ programs and 
%% Given that this work explores LLMs for performance prediction, our task is adjacent to the task of \textit{execution simulation}, where LLMs are used to predict and keep track of the internal state of a program.
%% These works are mostly Python-based, 
%% CoSm
%% REval
%% LiveCodeBench
%% PredEx
%
%
%
%
%\textbf{Shortcomings of Existing Work}
%
%\begin{enumerate}
%    \item \textbf{Limited Test Sets}
%    \item \textbf{Roofline Modeling Not Explicitly Considered}
%\end{enumerate}
%
%
%\textbf{The Roofline Performance Model.}


\section{Related Work}
\label{sec:RelatedWorks}

Our work lies at the intersection of Roofline-guided GPU performance modeling, execution-free prediction, and LLM-based code reasoning.
The closest prior work studies performance prediction without full profiler access, while adjacent areas use LLMs for HPC code generation, iterative optimization, and program-behavior reasoning.

%\textbf{Roofline modeling and static GPU performance prediction.}
%The Roofline model is a standard framework for reasoning about whether a kernel is limited by computation or data movement, and arithmetic intensity is one of its central quantities \cite{samuelwilliamsRooflineInsightfulVisual2009}. 
%Before LLMs, prior work explored predicting GPU performance without executing the target kernel, including static PTX-based execution-time estimation and learned models over PTX for performance, power, and energy prediction \cite{gargialavaniPredictingExecutionTime2018,guerreiroGPUStaticModeling2019}.
%More recent learned predictors extend this idea to unseen workloads and devices through few-shot transfer, hybrid GNN/LLM models, and deep-learning performance forecasters \cite{deyRelativePerformancePrediction2024,leeForecastingGPUPerformance2025,selvamCanLLMsEnhance2024,pannerselvamPerformancePredictionModels2025}.
%These works establish the value of execution-free prediction, but they do not ask whether off-the-shelf (not-finetuned) LLMs can estimate a quantitative Roofline quantity directly from GPU kernel source.
\textbf{Roofline modeling and static GPU performance prediction.}
The Roofline model is a standard framework for reasoning about whether a kernel is limited by computation or data movement, and arithmetic intensity is one of its central quantities \cite{samuelwilliamsRooflineInsightfulVisual2009}.
Long before LLMs, prior work studied GPU performance prediction using analytical and learned models built from program structure, microarchitectural assumptions, or low-level code features.
Early analytical models estimated execution time from memory- and thread-level parallelism or compiler-derived workflow graphs that capture effects such as divergence, coalescing, and bank conflicts \cite{hongAnalyticalModelGPU2009,baghsorkhiAdaptivePerformanceModeling2010}, while later static approaches predicted execution time directly from PTX or other compile-time features \cite{limAutotuningGPUKernels2017,alavaniPredictingExecutionTime2018,guerreiroGPUStaticModeling2019a}.
Roofline-style upper-bound analysis has also been extended to GPUs through cache-aware performance, power, and energy models \cite{lopesExploringGPUPerformance2017}, and newer predictors combine PTX features with learned models such as random forests or XGBoost, or replace execution-based profiling with static estimation of execution statistics \cite{braunSimpleModelPortable2021,emondsCUDAsapStaticallyDeterminedExecution2023,tranAccurateLatencyPredictor2025}.
Related work also explores profiling-informed analytical or machine-learning predictors, including bulk-synchronous-based models, DVFS-aware predictors, and more detailed analytical models of modern GPU cores \cite{amarisComparisonGPUExecution2016,wangGPGPUPerformanceEstimation2020,leeGCoMDetailedGPU2022}.
More recent learned predictors extend this direction to unseen workloads and devices through few-shot transfer, hybrid GNN/LLM models, and deep-learning performance forecasters \cite{deyRelativePerformancePrediction2024,leeForecastingGPUPerformance2025,selvamCanLLMsEnhance2024,pannerselvamPerformancePredictionModels2025}.
These works establish the value of execution-free prediction, but they mostly build specialized predictors for runtime, power, or energy from handcrafted models or trained feature pipelines.
Our work instead asks whether off-the-shelf LLMs can, through prompting alone, recover a quantitative Roofline signal of Arithmetic Intensity that is useful for early performance triage before profiler feedback is available.

\textbf{LLMs for HPC and parallel code understanding/generation.}
A broad line of work studies whether LLMs can better represent and generate HPC and parallel code through domain specialization, tailored tokenization, and benchmark-driven evaluation.
Early domain-specific efforts include \textit{Scope Is All You Need} and \textit{MonoCoder}, while HPC-Coder and related work show that HPC-focused corpora and fine-tuning can improve performance on parallel-code tasks \cite{talkadoshScopeAllYou2023,talkadoshMonoCoderDomainSpecificCode2023,danielnicholsModelingParallelPrograms2023,nicholsHPCCoderModelingParallel2024a}.
Other work studies compiler-oriented representations or benchmark suites for parallel and accelerator code, including Meta LLM Compiler, ParEval, KernelBench, and MultiKernelBench \cite{cumminsMetaLargeLanguage2024a,nicholsCanLargeLanguage2024,ouyangKernelBenchCanLLMs2025,wenMultiKernelBenchMultiPlatformBenchmark2025}.
Surveys and perspective papers likewise argue that HPC is a distinct setting for LLMs because correctness, hardware sensitivity, and performance portability all matter simultaneously \cite{chenLandscapeChallengesHPC2024,gongLanguageModelsCode2025}.
Relative to this literature, our goal is not to generate kernels, but to infer an interpretable performance signal from existing kernels before optimization begins.

\textbf{LLMs for iterative GPU code optimization.}
A large neighboring literature uses LLMs to improve code performance through iterative search with compilation, execution, and profiler feedback in the loop.
This includes evaluation-oriented studies such as KernelBench \cite{ouyangKernelBenchCanLLMs2025}, \texttt{robust-kbench} \cite{langeRobustAgenticCUDA2025}, and broader assessments of LLMs on HPC optimization tasks \cite{bowencuiLargeLanguageModels2025,cuiComprehensiveEvaluationLLMs2025}, as well as feedback-driven systems such as CUDA-LLM \cite{chenCUDALLMLLMsCan2025}, PerfCodeGen \cite{pengPerfCodeGenImprovingPerformance2025}, SpeedGen \cite{purschkeSpeedGenEnhancingCode2025}, WarpDrive \cite{damaniWarpDriveAgenticWorkflow}, IntelliPerf \cite{awadIntelliPerfLLMPoweredAutonomous2025}, Autocomp \cite{hongAutocompPowerfulPortable2025}, SwizzlePerf \cite{tschandSwizzlePerfHardwareAwareLLMs2025}, CudaForge \cite{zhangCudaForgeAgentFramework2025}, ParaCodex \cite{kaplanParaCodexProfilingGuidedAutonomous2026}, SysLLMatic \cite{pengSysLLMaticLargeLanguage2025}, the minimal-executable-program approach of Chu \textit{et al.} \cite{chuGPUKernelOptimization2025}, and agent-system approach in \cite{weiImprovingParallelProgram2024}.
A related training-based line uses RL or fine-tuning to optimize GPU code more effectively, including Kevin, CUDA-L1, KernelBlaster, and recent work on integrating performance tools into model reasoning \cite{baronioKevinMultiTurnRL2025,liCUDAL1ImprovingCUDA2026,dongKernelBlasterContinualCrossTask2026,nicholsIntegratingPerformanceTools2025}.
More broadly, evolutionary and autonomous search formulations such as LLM-GP and AlphaEvolve reinforce the value of iterative feedback for improving code quality and performance \cite{hembergEvolvingCodeLarge2024,novikovAlphaEvolveCodingAgent2025}.
Collectively, these works show that LLMs can optimize code when hardware feedback is available, whereas our work asks what useful guidance can be obtained before such feedback exists.
The downstream objective would be to later use LLM predictors in place of actual hardware profiling feedback to help drive these optimization tools.

\textbf{LLMs for GPU performance prediction.}
A smaller but much closer body of work uses LLMs as performance predictors rather than optimizers.
LLMPerf \cite{nguyen-nhatLLMPerfGPUPerformance2024} is an early demonstration that fine-tuned LLMs can act as GPU performance estimators for source programs.
Bolet et al. \cite{boletCanLargeLanguage2025} framed LLM-based prediction as a coarse Roofline-classification task that predicts whether kernels are compute-bound or bandwidth-bound.
More recently, Omniwise \cite{wangOmniwisePredictingGPU2025} introduced a fine-tuning pipeline for predicting several GPU kernel metrics, including bandwidth, cache behavior, GFLOPs, and arithmetic intensity, directly from code.
Our work differs by focusing on quantitative arithmetic-intensity prediction as an interpretable Roofline signal, by studying the zero-training regime with off-the-shelf models, and by explicitly analyzing the gap between source-level reasoning and compiler-realized behavior.
This distinction matters because floating-point semantics and lowered implementations of operations such as division can change effective FLOP counts in ways that are not explicit in source text \cite{IEEEStandardFloatingPoint2019,louvetNewtonRaphsonAlgorithmsFloatingpoint2010}.

\textbf{Program reasoning and code-property modeling.}
Arithmetic Intensity prediction also relates to whether LLMs can infer latent execution behavior and nontrivial code properties from source.
Prior work on code execution, predictive execution, and runtime-behavior reasoning shows that models can sometimes simulate program behavior, but often struggle as reasoning depth and trace complexity increase \cite{liuCodeExecutionPretrained2023,lyuLargeLanguageModels2024,liBlendedAnalysisPredictive2025,chenReasoningRuntimeBehavior2025}.
Benchmarks for static-analysis-style reasoning and broader code evaluation similarly show that current code LLMs remain brittle on deeper semantic tasks \cite{xieCOREBenchmarkingLLMs2025,jainLiveCodeBenchHolisticContamination2024}.
Related work further argues that text-only code LLMs often struggle to model latent code properties such as performance, motivating structured or lower-level representations when reasoning about behavior that is not directly visible in source \cite{pydaShortcomingsCodeLLMs2025}.
Our task fits naturally into this landscape because arithmetic-intensity prediction requires recovering execution-relevant behavior from source in a form that is directly useful for performance engineering.

\textbf{Gap addressed by this work.}
Taken together, prior work shows that static and learned models can predict aspects of GPU performance without execution, that LLM-based optimization systems benefit strongly from profiler feedback, and that LLMs remain imperfect at reasoning about latent program behavior.
What remains underexplored is whether off-the-shelf LLMs can serve as \emph{hardware-free Roofline triage tools} by estimating a quantitative profiler-relevant metric directly from code well enough to guide where scarce profiling and tuning effort should be spent.
Our work addresses this gap by studying Roofline Arithmetic Intensity prediction across GPU architectures, highlighting code features that make prediction harder, and identifying when post-compilation artifacts are needed to close the source-to-performance gap.

